\documentclass[10pt]{article}
\usepackage[T1]{fontenc}
\usepackage[margin=0.8in]{geometry}
\pagestyle{empty}
\setlength{\parindent}{0pt}
\setlength{\parskip}{2pt}
\begin{document}

\textbf{Qiong Yan; Peide Yang (corresponding author)}\par
Fujian Jiangxia University, Fuzhou 350108, China; Xiamen University Tan Kah Kee College, Xiamen, China\par
Emails: 381007382@qq.com (Qiong Yan); yangpd@xujc.com (Peide Yang, corresponding author)\par
\vspace{2pt}
08 November 2026\par
\vspace{2pt}
\textbf{To the Principal Editor, Computer Physics Communications}\par
\vspace{2pt}
Dear Editor,\par
\vspace{2pt}
Please consider the enclosed Computer Programs in Physics submission ``A differentiable, Poynting-consistent transfer-matrix and drift--diffusion kernel for co-design of thin-film solar cells'', by Qiong Yan and Peide Yang, for publication in Computer Physics Communications. The work is original, has not been published previously and is not under consideration elsewhere.\par
\vspace{2pt}
The manuscript describes UniDD, a deterministic Python program that joins a fully coherent, Poynting-consistent transfer-matrix optical solver to a self-consistent one-dimensional drift-diffusion solver in a single reverse-mode automatic-differentiation graph. This journal published the end-to-end differentiable photovoltaic simulator of Mann et al. (CPC 272, 108232, 2022), which uses Beer-Lambert absorption; UniDD extends that line of work by retaining thin-film interference inside the differentiable composition, without a surrogate and without interpolating the optical field onto the electrical mesh.\par
\vspace{2pt}
The numerical contribution is the solver construction and its validation: log-density variables, exponentially factored Scharfetter-Gummel currents, a fully coupled equilibrated and damped Newton iteration converged on a physical residual scale, and an implicit-function-theorem adjoint implemented as a custom torch.autograd.Function. External checks give an optical energy balance of 1.44e-6, a Beer-Lambert limit of 1.57e-8, composite-gradient agreement to 2.78e-3 and an adjoint residual of about 3e-14, with a gradient evaluation costing 1.10 times a forward solve (about 10\% overhead). The complete suite is deterministic at seed 0 and reproduces in about 25 minutes on one CPU core.\par
\vspace{2pt}
For reproducibility, the source archive includes the recorded optimisation histories together with the negative result: the pre-specified tenfold evaluation-count target is not met in eight dimensions (about threefold at moderate accuracy, neither optimiser reaching 1e-4). All targets are synthetic inverse crimes, and no experimental calibration is claimed. A SCAPS-1D 3.3.12 comparison was executed; two of four pre-specified criteria pass, and the illuminated J-V curves are not comparable because the contact models differ. The Program Summary, source code, result files, figure scripts and consistency checker accompany the submission under the MIT licence; the repository at https://github.com/qiantanzhijiao/TMM-DD-Codesign is made public at the time of submission.\par
\vspace{2pt}
A one-factor 8-D/16-D control (E4-C) and a noisy multi-start comparison against five derivative-free baselines (E5) are reported as independent supplements. They are shipped with their recorded histories and do not revise the negative E4 verdict.\par
\vspace{2pt}
Both authors approved the submission. The work is original, has no competing interests and received no specific funding. Generative AI was used only for language editing and manuscript source preparation, as disclosed in the manuscript.\par
\vspace{2pt}
Enclosures: manuscript (LaTeX and PDF), Program Summary, 5 figures, 6 tables, Highlights, and the source archive with README, MIT licence, environment files and a complete sample input/output test case; declarations are in the manuscript.\par
\vspace{2pt}
Yours sincerely,\par
\vspace{2pt}
\textbf{Peide Yang (corresponding author)}\par
Xiamen University Tan Kah Kee College; yangpd@xujc.com\par

\end{document}
